% ---------------------------------------------------------------------
%  Chapitre 6 — BLOC FINAL (esquissé)
% ---------------------------------------------------------------------
\chapter{Bloc final --- régression sur les signaux validés}
\label{chap:blocfinal}

% [À DÉVELOPPER] Principe : ne combiner QUE les signaux validés par leur
% bloc. Méthodes vues en cours mobilisables : régression multiple,
% régression logistique (cible win 0/1), sélection de variables (AIC,
% VIF, stepwise), validation respectant l'ordre temporel.

\section{Principe}
La régression finale ne combine que les signaux \emph{déjà validés} par
leur bloc respectif, afin de n'introduire aucun bruit dans le modèle
prédictif.

\section{Méthodes envisagées}
\begin{itemize}
  \item \textbf{Régression logistique} si la cible retenue est
        \texttt{win} (binaire)~: coefficients interprétables en
        \emph{log-odds}, sélection par AIC et contrôle de la colinéarité
        (VIF).
  \item \textbf{Régression pénalisée} (Lasso) pour la sélection de
        variables en présence de nombreux prédicteurs corrélés.
\end{itemize}

\begin{conditions}
En contexte temporel, la validation croisée doit \textbf{respecter
l'ordre chronologique} (pas de fuite d'information du futur).
\end{conditions}

\section{Résultats}
\emph{[À compléter une fois les blocs amont finalisés.]}
